\section{Supporting robustness and boundary results}
\label{app:supporting-results}

This section collects results that sharpen the main mechanisms without
being needed to reconstruct the paper's central argument. Their formal
counterparts and axiom dependencies are recorded in the formal-verification
supplement.

\subsection{Parametric robustness of the corner equilibrium}

\begin{proposition}[Perpetual sanctions on a patience interval]
\label{prop:perpetual-interval}
There is a non-degenerate interval of sender discount factors
$\beta_S\in[6/13,\,4/7)$ such that, for every $\beta_S$ in it, the
corresponding economy has a deterministic affine transition and admits a
stationary feedback Stackelberg equilibrium along which, from every initial state
of the high basin and at every date, the realized sanction stays above a
strictly positive floor, never falls, and the political state remains in
the basin. In fact, the realized sanction is maximal at every date.
\end{proposition}

This interval varies sender patience while holding the other primitives
fixed. It establishes robustness in that direction, not an open
neighborhood in the full parameter space.

\subsubsection{Symbolic corner economies}
\label{sec:existence:symbolic-family}

\begin{proposition}[Symbolically parameterized certified corner economies]
\label{prop:symbolic-family}
There is a class of economies whose sender and target stage payoffs carry
nine free real coefficients, with the adjustment rate, both discount
factors, and the basin geometry free as well, such that:
\begin{enumerate}[label=(\roman*)]
\item membership is defined by finitely many inequalities among these
primitives, some strict and some weak, and the region they define is non-empty;
\item every economy in that region carries a complete compliance certificate
in the legacy form and admits a stationary feedback Stackelberg equilibrium lying in
invariant sets of value functions whose radius and slope are closed forms in the primitives,
$K_S=(r+w)/\bigl(1-\beta_S(1-\delta)\bigr)$ and
$g_T=(n+v)/\bigl(1-\beta_T(1-\delta)\bigr)$; an economy with a strictly
interior basin bottom carries the
Milgrom--Shannon/general-$\F$ certificate as well; and
\item along that equilibrium the independently optimized myopic
equilibrium realizes no sanction at any feasible state, while the
forward-looking equilibrium realizes the maximal sanction at every
feasible state.
\end{enumerate}
\end{proposition}

\begin{proof}
Three conditions pin the equilibrium. An implacability condition makes the
forward-looking sender's dynamic objective strictly increasing in the
sanction at every feasible pair, so its response is maximal. A stage
condition keeps the slope of the sender's stage payoff in the sanction on the basin's low-sanction range
above the per-unit continuation drag. A entrenchment condition covers the target's
myopic sanction loss by its discounted continuation gain, making full
non-compliance the target's derived Bellman argmax. Because both policies are
at corners, the cross-Lipschitz coefficients vanish and the row
conditions reduce to $\beta_T<1$ and $\beta_S<1$. The Bellman images are
affine in the state, and the displayed constants of the invariant sets solve their
fixed-point equations. Non-emptiness follows by choosing each
coefficient strictly between bounds constructed from coefficients
chosen earlier; the formal supplement records the corresponding
Fourier--Motzkin elimination.
\end{proof}

The inequalities of (i) are of both kinds. The implacability and sender stage
conditions are strict, and so are $0<\delta<1$, $\beta_S,\beta_T<1$,
$q,u_0,\ell>0$ and $0<\bar s<1$ for the basin's sanction floor $\bar s$; the
conditions $\beta_S,\beta_T\ge0$, $q\le r$, $u_0+r\le k$, $n\ge0$,
$\ell\le c-n$, $0\le\bar x\le\bar s$ and the entrenchment condition are weak
(\path{SymbolicEconomy}).

These economies are corner-pinned. They do not contain the interior-sanction
benchmark or the dual-basin construction of \Cref{app:dual-basin-economy}.

The corner economies of \Cref{prop:primitive-persistence} contain
these: each symbolic economy, with the state slope bounds
$g_{\mathrm{lo}}=g_{\mathrm{hi}}=1-\delta$ and the lower bound on the sanction slope
$d_{\mathrm{lo}}=\delta$, is one of them
(\path{SymbolicEconomy.toF}). Clauses (ii) and (iii) hold on all the corner economies
as follows (\path{SymbolicEconomyF.symbolic_family_generalF}). At every corner economy
an equilibrium lies in invariant sets with the closed-form constants
$K_S=(r+w)/(1-\beta_Sg_{\mathrm{hi}})$, $g_S=(r+w)/(1-\beta_Sg_{\mathrm{lo}})$
and $g_T=(n+v)/(1-\beta_Tg_{\mathrm{hi}})$, and it realizes the maximal sanction
at every feasible state while the economy's own optimized myopic equilibrium
realizes none, with no side condition (\path{regime_separation_symMyopic}); a
corner economy with a strictly interior basin bottom carries the
Milgrom--Shannon/general-$\F$ certificate. Against the canonical myopic
equilibrium the separation is proved only under a strictly interior basin bottom
and the strict substitution condition $u_0+r<k$, which it reads at the top state. The
legacy-form certificate of (ii) does not extend, because its consequences use
the partial-adjustment law, and (i) is not claimed for the corner economies as a whole.

\begin{proposition}[Strict escalation]
\label{prop:strict-escalation}
Fix any economy, any equilibrium, and any interval $[a,b)$ on which the
equilibrium closed loop is invariant, and suppose that at every state of
$[a,b)$ a feasible sanction above the state strictly raises it (local upward
adjustment; it holds whenever the transition is strictly increasing in the
sanction, or satisfies \textup{(F5)} with $a\ge0$). If the equilibrium sanction
strictly exceeds the state throughout $[a,b)$, the political state
strictly rises at every date. If, in addition, the sanction rule is
strictly increasing on $[a,b)$, the realized sanction strictly rises at
every date, from every start in $[a,b)$, forever.
\end{proposition}

The two conditions are distinct: a rule can lie above the diagonal
without being monotone, or be strictly monotone while lying below it.
Because sanctions are bounded, strict escalation here means a strictly
increasing convergent sequence, not unbounded growth.

\subsection{Equilibrium feedback from sender patience}

\begin{proposition}[Equilibrium value feedback from sender patience]
\label{prop:equilibrium-patience-feedback}
Let $E_1$ and $E_2$ be two of the corner economies of
\Cref{prop:primitive-persistence}, which contain the symbolic economies of
\Cref{prop:symbolic-family}, that share the transition $\F$ and the sender's
stage payoff at the corner, $y\mapsto\Us(y,0,1)$. Write $\beta_1,\beta_2$ for
their sender discount factors and $V_S^{\beta_i}$ for the sender value of any
equilibrium of $E_i$ whose sender value lies in the economy's certified invariant set;
such an equilibrium exists at every corner economy (\Cref{prop:primitive-persistence}).
Along both equilibria the target policy and both realized actions remain
at their corners. Then, for every state $x$, writing $\F(x,1)$ for the
corner-pinned successor,
\begin{align}
V_S^{\beta_2}(x)-V_S^{\beta_1}(x)
={}&(\beta_2-\beta_1)\,
V_S^{\beta_2}\bigl(\F(x,1)\bigr) \label{eq:equilibrium-feedback}\\
&+\beta_1\left[
V_S^{\beta_2}\bigl(\F(x,1)\bigr)
-V_S^{\beta_1}\bigl(\F(x,1)\bigr)\right].\nonumber
\end{align}
At the absorbing lock-in state,
\begin{equation}
(1-\beta_1)
\left[V_S^{\beta_2}(1)-V_S^{\beta_1}(1)\right]
=
(\beta_2-\beta_1)V_S^{\beta_2}(1),
\label{eq:equilibrium-lockin-amplification}
\end{equation}
and $V_S^{\beta}(1)=U_S(1,0,1)/(1-\beta)$;
\eqref{eq:equilibrium-lockin-amplification} and the sign comparison that follows
need only a common corner payoff $U_S(1,0,1)$, not a common transition.
Consequently, if
$\beta_1<\beta_2$, then
$V_S^{\beta_1}(1)<V_S^{\beta_2}(1)$ if and only if
$U_S(1,0,1)>0$. The latter inequality holds at every corner economy,
so the strict ranking is unconditional on these economies. Finally, the
target equilibrium value is invariant to sender patience: if $E_1$ and
$E_2$ share $\F$, $\beta_T$ and the target's stage payoff at the corner,
$y\mapsto U_T(y,0,1)=-c+(n+v)y$, whatever their compliance costs $\ell$ and
their sender payoffs, then $V_T^{(1)}=V_T^{(2)}$ for any two such
equilibria.
\end{proposition}

Two remarks on the scope of the first two displays, because it is wider than the
set of economies they are stated over.

\emph{No functional form for $\F$ is used.}  These economies carry a general
transition, and \eqref{eq:equilibrium-feedback} and
\eqref{eq:equilibrium-lockin-amplification} require only that both value
functions solve the stationary Bellman equation along a common constant action
pair, against a common stage payoff and, for the first, a common transition. The
absorbing state that \eqref{eq:equilibrium-lockin-amplification} needs is
supplied by the diagonal condition $\F(1,1)=1$ rather than assumed, so it too is
free of functional form. Both are verified at that generality, and the versions for the corner
economies are special cases. On the symbolic economies of \Cref{prop:symbolic-family}, whose transition is affine,
the identities are also proved for their closed-form values, with the shared
adjustment rate in place of the shared transition.

\emph{The sign clause is an equivalence, and the direction it fails in is
economically real.}  Greater sender patience raises the lock-in value \emph{if
and only if} the sender's corner flow is positive. Where a locked-in sanctions regime pays
the sender a negative flow, patience makes the sender strictly worse off there,
and no feature of the dynamics reverses that. What makes the ranking
unconditional here is a property of the corner economies, their implacability condition, which
forces the corner flow to be positive; it is not a property of the mechanism.

\begin{proof}
At an equilibrium whose sender value lies in the certified invariant set the sender
sanctions fully at every state and the target maintains full non-compliance, since compliance
costs it $\ell>0$ per unit and does not move the sender off the corner; so both
values solve corner Bellman equations along the action pair $(a,s)=(0,1)$.
Subtracting the two sender equations gives \eqref{eq:equilibrium-feedback}. At
$x=1$, the successor state is one,
so the Bellman equation gives
$V_S^{\beta}(1)=U_S(1,0,1)+\beta V_S^{\beta}(1)$ and hence
\eqref{eq:equilibrium-lockin-amplification}. The sign comparison follows
from this identity. Of the conditions on these economies only implacability is needed for
$U_S(1,0,1)>0$: with the slope bound $g_{\mathrm{lo}}+d_{\mathrm{lo}}\le1$ it
places the worst stage disincentive $q+k-u_0$, which the sign condition
$u_0+r\le k$ puts at or above $q+r$, below the corner state slope $r+w$, so
$U_S(1,0,1)=u_0-q/2+r+w>u_0+q/2+r>0$. The target value is
invariant because its corner Bellman equation does not contain $\beta_S$ and has
exactly one bounded solution. Exact parameter
intervals are reported in \Cref{tab:compstat-windows}.
\end{proof}

\subsection{The upper-boundary trade-off}
\label{sec:absorbing-obstruction}

\begin{proposition}[Ratchet dichotomy]
\label{prop:ratchet-dichotomy}
Fix an economy in the model class, a sanction rule $R$ taking
values in $[0,1]$, and any interval $[\bar x,1]$ with $\bar x\le1$.
Exactly one of the following holds.
\begin{enumerate}[label=(\roman*),nosep,leftmargin=2.2em]
\item $R(1)<1$: the top sanction is interior. Then the ratchet floor
$R(x)\ge\max\{x,\bar s\}$ fails on $[\bar x,1]$, for every $\bar s$.
\item $R(1)=1$. Then $\F(1,R(1))=1$: the realized successor at the top
state is the top state itself, whatever $R$ does elsewhere.
\end{enumerate}
\end{proposition}

\begin{proof}
In the first case, evaluating the floor at $x=1$ would require
$1\le R(1)$, a contradiction. In the second, the diagonal condition
gives $\F(1,R(1))=\F(1,1)=1$.
\end{proof}

\begin{corollary}[Endpoint obstruction to an upper-basin floor]
\label{cor:endpoint-obstruction}
Let $R$ be any rule with $R(1)<1$. Then no interval
$[\bar x,1]$ can satisfy
\[
\max\{x,\underline s\}\le R(x)
\qquad\text{for every }x\in[\bar x,1].
\]
Consequently, any equilibrium whose realized sanction is zero at state
$1$ fails the dynamic-floor ingredient of the upper-basin persistence
certificate.
\end{corollary}

\begin{proposition}[Absorbing-boundary obstruction]
\label{prop:absorbing-obstruction}
Fix an economy in the model class and a candidate equilibrium. Suppose
the sender's continuation value rises to the top state at rate at least
$m$; shading the sanction below the corner by $1-s$ lowers the successor
state by at least $\kappa(1-s)$; and the sender's current-period gain
from that shading is at most $c(1-s)$. If
\[
c\le\beta_Sm\kappa,
\]
then the sender's response is pinned at $s=1$ across a
right-neighborhood of the target's action at $x=1$. Consequently, if
the target's stage payoff strictly increases when its action rises from
$\alpha(1)$ against the corner sanction, no such equilibrium exists.
\end{proposition}

At a pinned top sanction, the successor state is identical across the
relevant target actions. The target's continuation term therefore
cancels, leaving only its stage-payoff comparison. The proposition's
inequality supplies sufficient conditions for that pinning.

\subsection{A sufficient route to the weak switch leg}
\label{app:switch-route}

\Cref{app:two-channel} records that the sender-response-switch component is not
signed by the primitives. That stands. What follows is a \emph{sufficient} route,
one way of supplying the weak leg from conditions a modeller can check, and
not a derivation of the component's sign from the baseline.

Suppose that, across any two feasible sanctions, the \emph{difference} in the
sender's stage payoff is a function of the displacement from the sender's own
myopic best response alone, so that compliance relocates the sender's stage
problem without reshaping it, while an arbitrary compliance-dependent level is
left free and cancels from every sanction comparison. Write
$\Gamma_x(s)=\beta_S\Vs\bigl(\F(x,s)\bigr)$ for the discounted continuation term and
$\Lambda(a)=\sigDynTop_{\Vs}(x,a)-\sigM(x,a)$ for the response gap. Say that
$\Gamma_x$ has \emph{increasing translated chords} when the increment of $\Gamma_x$ across
a chord of any fixed length does not fall as the chord is translated toward
larger sanctions. Then, provided the shift-translated responses remain feasible
and the dynamic best response is taken as the greatest maximizer,
\[
\Gamma_x\ \text{has increasing translated chords}
\quad\Longrightarrow\quad
\Lambda\ \text{does not widen in compliance,}
\]
and the reverse chord condition makes $\Lambda$ not narrow. An affine $\Gamma_x$
satisfies both, leaving $\Lambda$ exactly constant. The proof is an exchange
argument (translate each maximizer by the shift and offer it in the other
problem) and uses no derivative, no strong concavity, and no compactness.
Convexity of $\Gamma_x$ is one way to supply the chord condition, and on an interval
with $\Gamma_x$ continuous it is equivalent to it; the chord form is preferred because
it names exactly the inequality the argument consumes.

The condition can be weakened further, and ordinally. Re-centring the sender's
problem on its own myopic best response, so that the gap \emph{is} the maximizer in
the displacement variable, the comparison becomes a monotone-comparative-statics
question, and what it requires is single-crossing lower dominance of the
re-centered objective across the two compliance levels. Increasing chords imply
that; the converse fails, since single crossing survives an arbitrary strictly
increasing transformation and a difference condition does not. This is the same
distinction \citet{DekelQuahSinander2025} draw in a model of costly adjustment:
monotonicity-type conditions govern the \emph{direction} of adjustment, and
curvature governs its \emph{speed}. The comparison here is a direction result.

\paragraph{Only the chords the decision problem generates are needed.}
The chord condition above is stated over every chord of $[0,1]$, and for a
continuous $\Gamma_x$ on an interval that is convexity, so as written it is not a
weakening of convexity at all. \citet{TaborTabor2009} explain why in general:
every nontrivial linear inequality valid for all convex functions is valid
\emph{only} for convex functions, once it is imposed over a family of affine
reparametrizations.

The exchange argument reads much less than that. It uses chords based at the two
myopic best responses $\sigM(x,a_L)$ and $\sigM(x,a_H)$ only, and only shifts inside
the interval those two compliance levels generate. Substituting
$u=\sigM(x,a_H)+z$ and writing $\Delta=\sigM(x,a_L)-\sigM(x,a_H)$ for the
sanction gap, that restricted condition says exactly
\begin{equation}
u\;\longmapsto\; \Gamma_x(u+\Delta)-\Gamma_x(u)
\quad\text{is non-decreasing on } [0,1-\Delta].
\label{eq:generated-chord}
\end{equation}
A \emph{single} lag, and that lag is the gap the two compliance levels
themselves generate. The all-chord condition asks the same of every lag. So what
buys the weakening is the restriction from all lags to one, and not a restriction
of the reparametrization group.

The distinction is not decorative, because it changes what a separating example
must look like. Tabor and Tabor's own escape ($x_1^2-x_2^2$, which satisfies
the spherical sub-mean-value inequality for every affine \emph{similarity} yet is
not convex) cannot serve here: it lives in $\mathbb{R}^2$ and works because the
function is harmonic, whereas in one dimension the corresponding inequality is
midpoint convexity, and a continuous midpoint-convex function on an interval is
convex. \Cref{eq:generated-chord} gives the construction directly instead: a
continuation term periodic with period $\Delta$ has vanishing $\Delta$-difference,
so it satisfies the condition however far from convex it is.

The weakening is therefore strict, and machine-checked to be. Take a sender whose
stage loss is single-peaked at $3/4-a/2$, so its stage argmax is a clamp falling
in compliance and the gap at $a_L=0$, $a_H=1$ is exactly $\tfrac12$; a transition
that adjusts fully; and the $\tfrac12$-periodic continuation value
$\varphi(s)=\bigl|\,|s-\tfrac12|-\tfrac14\bigr|$. It satisfies
\Cref{eq:generated-chord} and fails the all-chord condition, at base points
$\tfrac14$ and $\tfrac12$ with chord length $\tfrac14$, where the increment falls
from $\tfrac14$ to $-\tfrac14$. So the implication from the all-chord form does
not reverse.

Per the scope discipline this appendix keeps elsewhere, the example refutes a
universal and licenses nothing further: it does not say which condition holds at
any calibrated economy, nor that the weaker form is ever the binding one.

Three limits travel with all of this. Each implication is one-way, and a gap that
fails to widen does not force any chord condition. The relevant shape is that of
the \emph{composite} $\Gamma_x$, not of $\Vs$ or of $\F(x,\cdot)$ separately:
curvature of one factor says nothing about the composite without a matching
monotonicity of the other. And the response gap is \emph{not} the switch component:
it is one of five inputs to the switch-gap comparison, the remaining four being
shape conditions on the target's complete sanction payoff over the reached
sanctions.

That the ordering can genuinely fail is settled by an example economy whose
transition is strictly concave in the sanction, non-affine at every $\delta$, and
carries genuine inheritance, and in which a fully forward-looking equilibrium and
an independently optimized myopic one satisfy
$\alpha^{MM}(x)=0<1=\alpha^{FF}(x)$ at every feasible state, with both Bellman
equations holding exactly. Two things should be said about it, and no more.
It licenses exactly one claim (that the baseline together with several natural
monotonicity and optimality properties does not logically determine the two-sided
ordering), and nothing about whether the intended ordering holds in typical or
calibrated economies. And what it verifies is equilibrium-level: its
stage-payoff state shifts are chosen so that the proposed values solve the
Bellman equations exactly, and no claim is made that it satisfies the
forms of \Cref{ass:concavity} and its companions that quantify over sets of value functions.

\subsection{Why the path comparison assumes a monotone sanction rule}
\label{app:closed-loop-monotone}

The route taken here orders the state paths of the two equilibria by asking the
forward-looking sanction rule $R^{FF}(x)=\sigDynTop_{\Vs}\bigl(x,\alpha^{FF}(x)\bigr)$
to be non-decreasing in the state. That is a sufficient condition, not the exact
one: what the induction actually consumes is a cross-order condition on the two
closed-loop \emph{state} maps $x\mapsto\F(x,R(x))$, which a monotone
rule discharges. The monotonicity is assumed rather than derived, and the reason
is a constraint the model imposes rather than a gap in the argument.

Write the composition informally as $R'(x)=\sigma_x+\sigma_a\,\alpha'(x)$, where
$\sigma_a$ is the sender's response to compliance. This model maintains
$\sigma_a\le0$ throughout (compliance placates), so the two effects pull
against each other unless the target's own policy is non-increasing in the state.
The baseline does not deliver that. A target with stage payoff $-(a-x)^2$, which
does not mention the sanction at all, has $\alpha(x)=x$ as its genuine myopic
stage argmax, strictly \emph{increasing}. What that closes is the \emph{sign-only}
route: the closed-loop property cannot be read off the direction of $\alpha$ alone.
A quantitative route survives, since the direct state effect may dominate the
adverse compliance--sanction cross term even where $\alpha$ rises, but it turns on a comparison of
magnitudes the baseline does not supply. The honest options are to assume the
closed-loop property or to discharge it from a condition on the \emph{composite}. The
formalization takes the second: the single-crossing condition is imposed on the
sender's objective with the target's policy already substituted in, which is
precisely the formulation that never separates the two effects and so never
walks into the composition trap.

\paragraph{The condition the induction actually consumes.}
Write $G_i(x)=\F\bigl(x,R_i(x)\bigr)$ for the two closed-loop state maps. Stepping
from a lower state $x$ to an upper state $y\ge x$, the induction needs one
inequality, $G_1(x)\le G_2(y)$, and nothing else. Say that $G_1$ is
\emph{cross-ordered} below $G_2$ on a region $B$ when
\begin{equation}
x\in B,\quad y\in[0,1],\quad x\le y
\quad\Longrightarrow\quad
G_1(x)\;\le\; G_2(y).
\label{eq:cross-step}
\end{equation}

\begin{proposition}[Stopped path order under cross-step domination]
\label{prop:cross-step-paths}
Fix two sanction rules mapping $[0,1]$ into itself, a region $B\subseteq[0,1]$, a
feasible initial state, and a date $n$. If \eqref{eq:cross-step} holds and the
lower path has in fact remained in $B$ through step $n$, then the lower state path
lies weakly below the upper one at $n$. If in addition the same cross-order
condition holds for the two \emph{rules} in place of the two closed loops, the
sanction levied along the lower path at $n$ is weakly below the one levied along
the upper path.
\end{proposition}

Neither rule is required to be monotone, no invariance of $B$ is asserted, and $B$
need not be closed or an interval. What replaces those is a hypothesis about the
trajectory that was actually run, at the date in question.

Monotonicity of the upper rule, which \Cref{app:closed-loop-monotone} uses, is one
way to supply \eqref{eq:cross-step}: it gives $G_1(x)\le G_2(x)\le G_2(y)$ through
the transition's monotonicity in each argument. It is a strictly stronger way.
Condition \eqref{eq:cross-step} is equivalent to $G_1$ being dominated on $B$ by
the greatest non-decreasing minorant of $G_2$, the pointwise infimum of $G_2$
over the states weakly above the current one, and that minorant coincides with
$G_2$ exactly when $G_2$ is already monotone. Where it does not, the condition is
correspondingly weaker.

That gap is inhabited. Under the lattice-meet transition $\F(x,s)=\min\{x,s\}$,
which satisfies \textup{(F1)}--\textup{(F4)} but not \textup{(F5)} (it is
insensitive to a sanction above the state: $\F(0,1)=0$), the rule
$R(x)=\max\{x,1-x\}$ is feasible and \emph{not} monotone (it falls from $1$ at
$x=0$ to $\tfrac12$ at $x=\tfrac12$), yet it never names a sanction below the
state, so the closed loop is the identity and \eqref{eq:cross-step} holds
trivially. The point is that rule monotonicity is a property of the rule, while
what the comparison needs is a property of the rule composed with the transition;
the two come apart wherever the transition is insensitive to the sanction.

The two path conclusions are stated separately on purpose. Ordering the states
along the two paths does not by itself order the sanctions levied along them,
because the two rules are evaluated at different states, which is why the
sanction claim carries its own hypothesis rather than being read off the first.

The counterexample refutes derivability. It does \emph{not} exhibit an economy in
which the closed-loop rule actually fails to be monotone; that would
also require the effect of compliance on the sender's return to sanctioning to overturn the direct state effect.
The general shape of the obstruction is familiar from monotone comparative
statics, where what a composition needs is the follower's response to the
\emph{leader's action} and not only to the state.

\subsection{A sender payoff rising with the state, and a weaker persistence floor}
\label{sec:dynamic-stake}

For a fixed target policy, call the greatest maximizer of the sender's
one-period problem at a state the \emph{stage ceiling}.

\begin{proposition}[Dynamic floors require a sender payoff that rises with the state]
\label{prop:dynamic-stake}
Fix an economy in the model class, a continuous target policy, and the
sender's continuation value determined by that policy. If the sender's
stage payoff along the policy is non-increasing in the target's political
state, then at every state the greatest maximizer of $J_S$
is at most the stage ceiling. Equivalently, if at some state the
equilibrium sanction strictly exceeds the stage ceiling, then there are
feasible states $y_1\le y_2$ and a feasible sanction $s$ with
\[
U_S\bigl(y_1,\alpha(y_1),s\bigr)
<U_S\bigl(y_2,\alpha(y_2),s\bigr).
\]
\end{proposition}

The conclusion concerns the composed payoff
$y\mapsto U_S(y,\alpha(y),s)$, not the partial effect of the state at
fixed compliance. It also locates the required increase somewhere in
the feasible state space, not necessarily at the state where the
dynamic sanction exceeds the stage ceiling.

\subsubsection{Persistence under a level floor}
\label{sec:weak-floor}

\begin{proposition}[Level-floor invariance and its cost]
\label{prop:level-floor}
Suppose $R(x)\ge\bar s$ on $[\bar x,1]$ with $\bar x\le\bar s$.
Then $[\bar x,1]$ is forward invariant under the closed-loop
map of $R$, and the statewise ordering of dynamic over myopic sanctions and
the finite-shock entry conclusion are unaffected. The level floor does
not, however, entail the ratchet step
$x\le\F\bigl(x,R(x)\bigr)$ and therefore does not entail that the
political state is non-decreasing: there are transitions and rules
satisfying the level floor along which the state falls.
\end{proposition}

Forward invariance follows from monotonicity and the diagonal condition:
\[
\bar x=\F(\bar x,\bar x)
\le \F(\bar x,\bar s)
\le \F(x,\bar s)
\le \F\bigl(x,R(x)\bigr).
\]
For the partial-adjustment law with $\delta=1/5$, $\bar x=\bar s=1/2$, and a rule
constant at $1/2$, the state falls from $1$ to $9/10$. Thus the level
floor preserves the basin but does not imply a monotone state path.

\subsection{Further exposition}
\label{app:moved-ad}

None of the passages below is a theorem statement
or a hypothesis of a main-text result; each qualifies, illustrates or
inventories one. Each is headed by what it is and where it stood.

\paragraph{The statement on the cube and the formalization on the line.} \emph{(From \Cref{sec:model}, after \textup{(F1)--(F5)}.)}

One difference between this statement and the formalization should be
recorded rather than left to inspection. States and sanctions are
economically meaningful only on $[0,1]$, and \textup{(F1)} and
\textup{(F5)} are accordingly stated there. The formalization, however,
maintains \textup{(F2)} and \textup{(F3)} \emph{globally} ($\F$ is
monotone in each argument on all of $\mathbb{R}$, not merely on the
cube) because a global monotonicity hypothesis is markedly easier to
carry through the order arguments than its restriction. Read as an
assumption on $\F$ it is strictly stronger than the displayed one.
It costs no generality: a transition that is continuous, coordinatewise
nondecreasing, and diagonal-fixing on the square extends to one with the
same three properties on all of $\R^2$ and the same values on the
square, by adding to $\F(\operatorname{clip}x,\operatorname{clip}s)$ the
two correction terms $\min\{0,\min\{x,s\}\}$ and
$\max\{0,\max\{x,s\}-1\}$, which vanish on the square and restore
the diagonal outside it. That extension is machine-checked
(\texttt{Foundations/TransitionExtension}); every difference-quotient condition stated
on the square is inherited verbatim, since the extension \emph{is}
$\F$ there. The payoff functions and the stage selector extend by
clipping their arguments. Both steps are assembled into a machine-checked
constructor (\texttt{Postulates/ParamsOfCube}): from payoffs continuous on
the cube, a selector feasible, continuous and argmax on the square, and a
transition satisfying \textup{(F1)--(F4)} on the square, it builds a formal
economy that agrees with the data on the cube and square and inherits
\textup{(F5)} from it. So the formal statements quantify over a class no
smaller than the displayed one, and every theorem stated over the formal
economy applies to every cube-defined economy, evaluated where that economy
is defined.
The point is recorded because an assumption surface that is wider in
the machine-checked statements than in the prose is the kind of gap a
reader is entitled to be told about rather than to find.

\paragraph{Continuous payoffs whose myopic best response jumps.}

The formal economy asks of the myopic sender only a feasible stage maximizer
$\sigM(x,a)$, which exists whenever the payoffs are continuous; it does not ask
that this choice vary continuously with the state and the compliance level.
That continuity (\path{Params.StaticSelectorContinuous}) is read only by the
results built on the myopic sender: the target's value against the myopic
sender, the target-only case, the canonical and
greatest myopic equilibria, and every comparison with the canonical myopic
equilibrium, such as the three-goals conclusion of the route of
\Cref{prop:flip-route}. Those results take it as a hypothesis, which every
economy constructed in the formalization satisfies, and the class
keeps it as a condition of its own, because its
results on the four cases compare with the myopic sender. Every other result
stated for the general game (in particular \Cref{thm:existence},
\Cref{thm:corner-existence}, \Cref{thm:persistence},
\Cref{cor:nontermination}, \Cref{thm:shock}, \Cref{thm:finite-episode},
\Cref{thm:relief-exit}, \Cref{thm:lowbasin}, \Cref{prop:longrun},
\Cref{prop:order-selection}, \Cref{prop:history-deviations} and
\Cref{prop:strict-dominance}) holds for every economy with continuous
payoffs that meets that result's own hypotheses; none of them reads this
continuity. That includes an economy whose sender payoff is linear in the sanction
with a slope whose sign depends on the state, so that the myopic best response
jumps between no sanction and the full sanction. No such economy is constructed
in the formalization, so this generality is stated, not yet illustrated by an example.

\paragraph{Why differentiability is stated on a neighborhood of the cube.} \emph{(From the footnote to the payoff domains in \Cref{sec:model}.)}

The Lean development states every quantitative bound it uses
(Lipschitz constants, lower state slopes, curvature moduli) as a
finite-difference inequality on the compact cube, and for the polynomial
instances derives those bounds from the coefficients. Where a
smooth route is used to \emph{read off} such a constant, it needs
the second partials to be continuous up to the boundary, which is
why the hypothesis is stated on a neighborhood of the cube:
continuity on $[0,1]^3$ together with $C^2$ on the open cube alone
does not bound the second partials ($x \mapsto x^{3/2}$ is $C^1$ on
$[0,1]$, $C^\infty$ on $(0,1)$, and has unbounded second
derivative). A reader who prefers a weaker maintained assumption may
drop the differentiability altogether and take the finite-difference bounds as
primitive; nothing else changes.

\paragraph{Which result consumes which structural condition.} \emph{(From \Cref{sec:assumptions}.)}

\begin{table}[htbp]
\centering\footnotesize
\caption{Which result consumes which structural condition.}
\label{tab:assumption-routes}
\begin{tabular}{@{}p{0.36\linewidth}p{0.60\linewidth}@{}}
\toprule
\textbf{Result} & \textbf{Conditions consumed}\\
\midrule
\Cref{thm:persistence} (persistence, general $\F$) & the package (P1)--(P5): membership of the value sets (P1), the sender Lipschitz ceiling (P2), A1 in its value-set form (P3), the target combined slope condition A2 as (P4), the sub-pivot slope condition (P5); no A3. Dropping (P4) keeps feasibility and optimality, conclusions (iv)--(v), from the equilibrium object and loses only the preference conclusion (vi); that P4-free statement is a proved composition (\Cref{app:formalization}), not a second signature. A1 is read only for (iii): (i), (ii) and (iv) follow from (P1), (P2) and (P5).\\
\Cref{thm:shock} (minimum-drift shock) & (P1)--(P3) and (P5), with no target condition, together with the transition's (F1)--(F4), full-sanction ascent below the basin (which gives $\bar x<1$), and the temporary shock; the package alone is not the entry theorem. No A3.\\
\Cref{thm:finite-episode}, \Cref{thm:relief-exit} & the equilibrium, transition, floor and rule hypotheses stated in each: (F1)--(F4), with held-sanction ascent below the basin for entry and post-basin descent for exit in place of (F5); no A1--A3 hypothesis appears in their statements.\\
\Cref{prop:longrun} (long run) & the level condition at an equilibrium (the sanction rule at or above the state on the basin, with no floor) plus (F1)--(F5). (P1), (P2) and (P5) imply the level condition; neither A1 nor (P4) is read.\\
\Cref{thm:existence} (Schur condition) & invariant sets of value functions; dynamic argmax uniqueness, the sender's at feasible compliance and the target's at every value pair; joint-operator invariance of those sets; the row bounds with the Schur condition \eqref{eq:Schur} (\Cref{app:existencemachinery}). A3 gives only stage-argmax uniqueness, neither the dynamic uniqueness assumed nor a substitute for the other premises.\\
\Cref{prop:flip-route} (derived-basin route) & the sender half of \Cref{ass:flip} on the floor range and the sender's top boundary sign of \eqref{eq:boundary} derive the sender pivot; the three-goals conclusion consumes in addition A1 on the derived basin, row-sum contraction ($\beta_T+L_{TS}<1$, $\beta_S+L_{ST}<1$), full-sanction ascent below the basin, and the diagonal-residual, monotone-sender-values and myopic-ceiling clauses (i)--(iii) (\Cref{app:dischargers}); no target preference or sign condition and no A3. The target combined slope condition A2 can still be derived on the derived basin from the target half and the target's top sign, but the route does not read it.\\
\Cref{prop:placable-family}, \Cref{prop:target-only} (cases $00$, $10$) & the class's fields of \Cref{sec:four-regimes}, the concavity condition and the Riccati closure; as hypotheses of the results that read them, the myopic relief condition, interiority where $A/q<1$ or a family closed form is read and elsewhere only myopic interiority $u_{\mathrm{hi}}<q$, the lifting condition for case $00$ (the myopic equilibrium itself reads only the myopic lifting condition $u_{\mathrm{hi}}<k_{\mathrm{lo}}$) and the entrenchment condition in place of A2 for case $10$.\\
\Cref{thm:regime-separation} (case $11$) & the same class fields, the myopic relief condition, myopic interiority $u_{\mathrm{hi}}<q$, the myopic lifting condition $u_{\mathrm{hi}}<k_{\mathrm{lo}}$, the entrenchment condition and a non-decreasing response at non-compliance, which the increasing-differences condition \eqref{eq:rally-margin} implies.\\
\Cref{prop:sender-only} (case $01$) & the same class fields with interiority, the lifting condition, the sender-only compliance condition \textup{(T6)}, the forward-looking relief condition \textup{(RL)} of \Cref{app:proofs:placable} and compliance neutrality.\\
\Cref{thm:four-cells} & the union of the three rows above. None of the four uses A1--A3 as stated or \Cref{ass:flip}: the class's signs $r_{\mathrm{lo}},n_{\mathrm{lo}}\ge0$ permit strict complementarity, and what it lacks is the boundary sign crossing.\\
\Cref{thm:dominance} & a sender value non-decreasing in the state; no equilibrium and no basin; none of A1--A3.\\
\Cref{prop:strict-dominance} & one comparison of the sender's continuation value at the successor state per lower sanction (no sanction at or below the greatest stage maximizer carries a larger continuation value at the successor state than it) together with one sanction at which the stage loss measured from the greatest stage maximizer is outweighed by the discounted continuation gain; the comparison alone can hold with equality and change no choice.\\
Interest-group polity (\Cref{app:fourgroup}) & derives the cross-partial signs of \Cref{ass:flip} and the curvature of A3 from region (R); A2 through the entrenchment-condition route at the certificate-complete instance, not from \Cref{ass:flip}. The welfare sign of \Cref{app:fourgroup} is separate: region (R) implies neither of the two welfare inequalities.\\
\bottomrule
\end{tabular}
\end{table}

\paragraph{The substitution between compliance and the sanction needs no concavity.} \emph{(From \Cref{sec:assumptions}, after A3.)}

A3 is not, however, what carries the substitution between compliance and the sanction. Decreasing differences of
$\Us$ in $(s,a)$ orders the stage argmax \emph{correspondence} in the strong set
order, and while that does not order an arbitrary selection from it, it does
order the \emph{greatest} maximizer, which exists by Weierstrass alone. So the
antitonicity of $\sigM$ in $a$ that the ordering results use follows from the
cross-partial sign plus a greatest-selection convention, with no concavity and no
uniqueness; uniqueness is merely one sufficient way to reach it, since a unique
maximizer is trivially the greatest. This matters because uniqueness is a
regularity condition rather than an economic primitive, and it fails exactly
where a sender payoff is linear in the sanction or flat over a range.

\paragraph{Static and dynamic thresholds, substitution, and the constructed pivots.} \emph{(From the end of \Cref{sec:flip}.)}

\paragraph{Static and dynamic sender thresholds.}
The construction at slope level $m=0$ gives the \emph{static flip}
$x_S^0:=\inf\mathcal F_S(\bar s;0)$: the state at which
sanctioning stops being a stage loss to the sender. Since a higher
slope level is a more demanding condition,
$x_S^0\le\bar x_S$:
between the static flip and the dynamic pivot, the sender's stage
incentive has already flipped but
the transition-adjusted forward-looking drag
\eqref{eq:F-continuation-allowance} is not yet dominated. The band is
indexed by the slope level $m_S$ and is non-decreasing in it, since a
higher level is a more demanding condition. Under the partial-adjustment law the
drag per unit sanction is $\beta_S K_S\delta$, which vanishes in the
myopic limit $\beta_S\to0$, where the band closes; no comparative
static in $\beta_S$ is proved beyond that observation.

\paragraph{Compliance placates.}
Suppose compliance lowers the sender's marginal value of sanctioning:
\[
\frac{\partial^2\Us}{\partial s\,\partial a}\le0.
\]
Define
$x_S(a)$ as the infimum of the states at which the sender's difference-quotient inequality in
\eqref{eq:FS-def} holds for that fixed action $a$, rather than uniformly
over all actions. The resulting per-action threshold
$a\mapsto x_S(a)$ is non-decreasing: against a more compliant target,
a harder state is required before sanctioning clears the margin. The
uniform flip set is therefore determined by the full-compliance slice,
so $\bar x_S=x_S(1)$.

An example economy has
$\bar x_T=2/5$, $\bar x_S=1/2$, and a target sanction slope that changes
sign at the interior pivot. The construction uses a kinked sender payoff
because its diagonal-residual and myopic-ceiling conditions cannot both
hold smoothly near the top state.

Because the pivots are infima of upper flip sets, their comparative
statics follow by set inclusion. A payoff change that raises the
target's sanction slopes on the high range lowers $\bar x_T$, and the sender
analogue shifts the interval between its static and dynamic thresholds.
The sanction floor $\bar s$ remains
a separate parameter in the construction; it may be set equal to the
derived basin pivot only after the pivots have been computed.

\paragraph{The four cases: arithmetic form, sanction rules, and the one-period ranking.} \emph{(From the paragraph after \Cref{thm:four-cells}.)}

Two arithmetic
consequences follow from the constructed rules alone. For the limit of the state,
writing $L^{ij}$ for the four limits from a common start,
$(L^{11}-L^{10})-(L^{01}-L^{00})=L^{11}-L^{10}>0$. At a common state $x$,
writing $R^{ij}(x)$
for the four sanction rules, $R^{00}=R^{01}=0$, $R^{10}(x)=\sigM(x,0)$
and $R^{11}(x)>\sigM(x,0)$ (the vertex $(u_0+rx)/q$ on the quadratic--linear
family), so
$(R^{11}-R^{10})(x)-(R^{01}-R^{00})(x)>0$. Both are restatements of the rule
identities, not interaction primitives. At every state the realized sanction is
zero in the two cases in which the target is myopic and buys
relief, equals the positive myopic sanction when only the target is forward-looking, and
lies strictly between the myopic sanction and one when both are, above it by
at least the foresight premium \eqref{eq:sender-wedge}; with
$a^{01}(x)>a^{00}(x)$, the sender's forward-looking incentive raises the sanction against a non-compliant
forward-looking target and extracts more compliance from a myopic one. \citet{Eaton1992} find that a more patient sender is more likely to succeed
in their threat model; what is compared here is a pair of constructed
equilibrium rules at a common state, not their toughness parameter or a threat
stage, so the two statements are neighbors and not the same result. The one-period ranking
$\ell_{\mathrm{hi}}<\underline s\,p_{\mathrm{lo}}$, under which full compliance beats full non-compliance in the one-period problem, is sufficient and not necessary for both forward-looking incentives to change behavior: on the
enriched subclass, the increasing-differences condition \eqref{eq:rally-margin} and the forward-looking relief condition \textup{(RL)} alone give it,
through a small compliance rather than full compliance, and \textup{(RL)}
does not imply the one-period ranking: the affine relief-only member satisfies
\textup{(RL)} as $6<12$ and fails the ranking as $6<4$.

\paragraph{Two boundaries of the sender-only case.} \emph{(From after \Cref{prop:sender-only}.)}

Two boundaries of the case are worth stating. First, \Cref{prop:sender-only}
gives $a^{00}(x)<a^{01}(x)\le1$ and no more: at an economy in the class with
the lifting condition at equality ($A=k$) and no state--sanction cross terms, on the quadratic--linear
family, the zero level is full
compliance, $a^{01}(x)=1$ at every state, while $a^{00}=3/4$
(\Cref{app:formalization}). Relief from a
forward-looking sender is bought at a strictly higher price, which can be the
whole of compliance. Second, the entrenchment condition that makes the
forward-looking target maintain non-compliance is sufficient, not necessary: with the same payoff
shapes and the target's state valuation lowered until the condition fails,
non-compliance is still the unique optimum at a constructed equilibrium
(\Cref{app:formalization}). Where the condition fails the
class's classification is silent, not reversed.

\paragraph{The Schur-condition existence route, summarized.} \emph{(From \Cref{sec:existence:general}.)}

\Cref{thm:existence} (\Cref{app:existencemachinery}) states the route: on
closed, non-empty invariant sets of value functions, which the joint Bellman operator maps into
themselves, with the sender argmax unique on the sender's invariant set at every
compliance level in $[0,1]$ and the target argmax unique at every value pair,
if the two Bellman rows are Lipschitz in
the value pair with coefficients $\beta_T,\beta_{S,\mathrm{row}}<1$ on the
diagonal and $L_{TS}$, $L_{ST}$ off it, all four of any sign, and the Schur
condition $L_{TS}L_{ST}<(1-\beta_T)(1-\beta_{S,\mathrm{row}})$ holds (for non-negative
coefficients, equivalently $\rho(\mathbf M)<1$), then the operator has a unique
fixed point on these sets in a weighted sup-norm, the associated optimal target
policy and the sender's dynamic best response form an equilibrium in $\mathcal E$ with
values in these sets, and every equilibrium with values in these sets that
plays the target selector those sets define coincides with it.

The Schur inequality permits one Bellman row sum to exceed one when the
opposing cross-value coupling is sufficiently weak. Its uniqueness
conclusion is confined to the stated invariant sets and selector restrictions;
it is not global equilibrium uniqueness. The appendix states the theorem and
gives the Perov reduction and the exact constants
(\Cref{app:existencemachinery}).

\paragraph{Off-corner existence and the transition-smoothed game.} \emph{(From \Cref{sec:existence:general}, after \Cref{thm:corner-existence}.)}

With an interior target action, each Bellman row depends on the other
player's continuation value. The deterministic joint iteration need not
preserve order and can cycle even when current argmaxes are unique. We
call the game with an atomless transition component the
\emph{transition-smoothed game} (``purified'' in earlier versions and in the
formalization's names); the term refers to transition smoothing, not payoff
purification.
Under the regularity conditions in \Cref{app:existencemachinery}, its
joint operator is a contraction and supplies an off-corner equilibrium
\citep{HeSun2017,DoraszelskiSatterthwaite2010}. Deterministic off-corner
existence is established for the example economies; the primitive
model by itself supplies no global existence theorem of that form.

\paragraph{The benchmark economy.} \emph{(From \Cref{sec:existence:nonvacuity}.)}

The first application shows that the conditions of \Cref{thm:persistence} are
attained in an economy where continuation values affect observable
choices: \Cref{thm:b2-benchmark} (\Cref{app:witness:benchmark}) exhibits,
within the general-$\F$ model, an economy with positive discount factors for
both players, a non-degenerate high basin and an equilibrium at which both
value functions are nonconstant, the target maintains full non-compliance at the basin pivot
while the sender's realized sanction there is strictly interior and strictly
above its myopic level, every basin start keeps the realized sanction above
the pivot and non-decreasing, and a finite temporary shock at the maximum
sanction carries every initial state into the basin. The theorem establishes
existence in the model class. The appendix
strengthens its robustness in two
directions: permanent sanctions survive on a non-degenerate interval of
sender discount factors, and the symbolic corner economies free the payoff
coefficients, adjustment rate, discount factors, and basin geometry
(\Cref{prop:perpetual-interval,prop:symbolic-family}).

\paragraph{The off-corner boundary of primitive persistence.} \emph{(From the end of \Cref{sec:existence:primitive-persistence}.)}

The off-corner constructions locate the remaining boundary. A purified
economy defined on the full state space has nonconstant values, a state-varying
target policy, and a monotone sanction rule, but its top-state
sanction is interior and therefore cannot satisfy a ratchet floor on a
basin ending at one. It nevertheless ratchets strictly on an interior
sub-basin. The supporting appendix states the general boundary results;
\Cref{sec:discussion:scope} explains their economic content. Thus
equilibrium existence and continuation values that change choices do not alone
imply the stronger conditions of \Cref{thm:persistence} on a basin ending at one.

\paragraph{A local form of the conditions of the high-basin theorem, and no selection of a limit.} \emph{(From after \Cref{thm:persistence}.)}

The same
conclusions (i)--(iii) also follow from three conditions stated at the
equilibrium pair alone (a Lipschitz modulus of the equilibrium sender
value, one comparison per sub-pivot sanction against the pivot at the
equilibrium policy, and interval dominance of the substituted objective)
with no value sets and no value class
(\path{standardCompliance_persistence_generalF_local_star}); P1--P3 and P5
supply those conditions at the pivot, and the local form constructs no
equilibrium. The first two conditions alone give conclusions (i) and (ii),
a non-decreasing state path and the target's optimality along it
(\path{standardCompliance_persistence_generalF_level_star}): interval
dominance is read only for (iii). The package does not select a limit:
\path{persistence_package_not_unique_steady_state} fills every field of
P1--P3 and P5 at an economy whose equilibrium closed loop is the identity on
the basin, so every basin state is steady and no state is a common limit of
the basin's paths; attraction is the separate content of
\Cref{prop:order-selection} and \Cref{prop:derived-interval}.

\paragraph{What the derived interval does not say.} \emph{(From after \Cref{prop:derived-interval}.)}

The result neither places $x_s$ in the interior nor describes starts above it. At
the example economy $x_s=1$, so the derived interval is the whole high basin
and the proposition specializes to the absorbing-corner case. The three results
therefore have distinct roles: \Cref{prop:longrun} proves convergence under the
level condition of the high-basin theorem, \Cref{prop:order-selection} selects an extremal limit,
and \Cref{prop:derived-interval} obtains a common limit on a derived interval.

\paragraph{Corners at the dual-basin economy.} \emph{(From after \Cref{prop:derived-dividing-line}.)}

At this example economy the target chooses full non-compliance and the realized
sanction lies at a corner on each basin, so both displayed long-run outcomes are
corners. These are features of the construction, not of the general selection
result, whose steady states may be interior. A sanction floor extending to the top
state forces the sanction rule to equal one at the top state, $R(1)=1$; that
the rule sits at the ceiling throughout the high basin is a feature of this
economy's payoffs, not a consequence of the floor (an equilibrium rule
$(x+1)/2$ satisfies such a floor with $R(x)<1$ below the top). This
absorbing corner is a specialization of the general result rather than the
state-dependent ratchet motivating the mechanism.

\paragraph{The pathwise ordering at the Eaton--Engers corner economy.} \emph{(From after \Cref{cor:pathwise-dominance}.)}

At the Eaton--Engers corner economy, the myopic
rule is
\begin{equation}
\sigM\bigl(x,\alphaM(x)\bigr)
=\max\Bigl\{0,x-\tfrac15\Bigr\},
\label{eq:ee-myopic-rule}
\end{equation}
so a high basin is not preserved. The pathwise ordering still holds there
for the simpler reason that the dynamic sanction is one and every
myopic sanction is feasible and hence no greater than one
(\path{eeCert_pathwise_dominance}). The proof of the corollary, which needs basin
preservation and so does not apply here, is in \Cref{app:proof-pathwise}.

\paragraph{The top of the state space, interior ratchets, and the open primitive problem.} \emph{(From \Cref{sec:discussion:scope}.)}

The transition structure also creates an economic trade-off at the top
of the state space. A ratchet floor on $[\bar x,1]$ requires the realized
sanction at state one to equal one. The diagonal condition then makes the
top state absorbing. If the sender response is pinned at the corner
across nearby target actions, those actions have the same successor
state, so the target's continuation term cancels from its choice. The
target policy at the boundary must then be justified by current payoffs.
\Cref{prop:ratchet-dichotomy,prop:absorbing-obstruction} state the
kinematic and incentive parts of this boundary.

An interior sanction at the top state is therefore incompatible with
that particular floor, but not with persistence in every form. A fixed
level floor preserves the high basin without guaranteeing a monotone
state path (\Cref{prop:level-floor}). The purified construction on the
full state space goes further: its political state and realized sanction
rise strictly on an interior sub-basin and converge to an interior steady
state. The ratchet cannot be extended to the ceiling because its
right-hand endpoint must be a fixed point. Off-corner ratcheting is thus
possible, but it terminates before the absorbing top boundary.

A strict ratchet is needed neither for the political state to settle nor
for the limit to be an extremal steady state; what it buys is a common
long-run outcome for a whole interval, which on an invariant interval is
equivalent to uniqueness of the steady state there. Deriving that
uniqueness from payoff and transition primitives on a pre-specified interval
is the open primitive problem; \Cref{app:scope-notes} records why the
corner proof cannot simply be relaxed and what a sufficient condition must
imply.

\paragraph{The North Korea mapping and what it leaves out.} \emph{(From \Cref{sec:discussion}, predictions and limitations.)}

The North Korea application is interpretive. One possible mapping is
that tighter control raises the influence of actors who benefit from
isolation (the model's political state), while resistance to opening
that would weaken such control corresponds to the target's lower compliance.
\citet{HaggardNoland2017} and \citet{UNPanelDPRK2024} document broader
political control and resistance to opening, not the model's state or the
sender's valuation of it directly. The same sources also document illicit
foreign-exchange channels, sanctions evasion, and international
enforcement and coordination problems, none of which is represented in
the one-state baseline. The model isolates one domestic political
feedback that can coexist with those mechanisms; it does not attribute
North Korea's outcomes to the political-state mechanism alone.

\paragraph{The theorem on a finite temporary shock: the comparison in (iii), two boundaries, and an instantiation.} \emph{(From after \Cref{thm:finite-episode}.)}

The comparison in (iii) is with the path under the sanction held at $h$: the
path under the shock need not be monotone (an early sanction can push the state
above a later one), but it dominates the held path at every date through
$T$, and the held path is monotone and gains at least $\eta_h(\bar x)$ per
period below $\bar x$. Two boundaries delimit the statement, both
machine-checked as refutations. A floor \emph{equal} to the pivot need not
enter in finite time: under the partial-adjustment law, a sanction
held at $\bar x>0$ from state $0$ stays strictly below $\bar x$ at every
date, so held-sanction ascent, which forces $\bar x<h$, cannot be weakened to
$h=\bar x$. And there is no
uniformly small duration: for every preassigned $T$ and pivot in $(0,1]$
there is a partial-adjustment law under which $T$ periods of full sanction from $0$ end
below the pivot. The theorem also says nothing about the dates before $T$,
at which the sanction is set exogenously and not chosen, and nothing about whether
the shock was worth imposing.

\paragraph{Instantiation.} At the example economy of the
general-transition certificate route with a non-affine transition
($\F(x,s)=x/2+s/2+(s-x)^2(1-s)/100$ on the square, pivot $1/10$), the
constructed equilibrium's sanction rule is at or above $1/10$ on
$[1/10,1]$, and (ii), the maximal-shock case of (iii), and the general
case of (iii) are all discharged (declarations in \Cref{app:formalization}).

\paragraph{Entry and exit durations compared.} \emph{(From after \Cref{thm:relief-exit}.)}

The entry
and exit durations are computed from different sections of the same
transition, $\F(\cdot,1)$ below the threshold and $\F(\cdot,0)$ above
it, so they need not coincide. Their difference is an asymmetry between
two sufficient duration bounds; it is a difference of durations, not a band
in the state space, and neither bound is a required duration. Under the
partial-adjustment law they collapse to
$\eta_\delta(\bar x)=\delta(1-\bar x)$ and $\eta^{R}_\delta(\bar x)=\delta\bar x$,
so before rounding the exit bound exceeds the entry bound when $\bar x<\tfrac12$
and falls below it when $\bar x>\tfrac12$: the asymmetry is governed by which
side of the midpoint the threshold sits on, not by $\delta$, which cancels. (As
in \Cref{cor:shock-transition-strength}, the comparison of the \emph{integer}
durations remains weak, since distinct bounds can share a ceiling; the
specialization to the partial-adjustment law is machine-checked as \path{affine_entryBound_lt_exitBound_iff},
with \path{affine_fullReliefMinDrift_eq} for $\eta^{R}_\delta$.) For a general
transition under strict adjustment and $0<\bar x<1$, the exit bound exceeds the
entry bound if and only if
$\bar x\,\eta^{R}_{\F}(\bar x)<(1-\bar x)\,\eta_{\F}(\bar x)$
(\path{entryBound_lt_exitBound_iff}). If the transition is symmetric under
reflection of the square, $\F(1-x,1-s)=1-\F(x,s)$, so that relief at $x$ mirrors
a shock at $1-x$, then $\eta^{R}_{\F}(\bar x)=\eta_{\F}(1-\bar x)$, the exit duration
at $\bar x$ is the entry duration at $1-\bar x$, and the exit bound exceeds the
entry bound exactly when $\bar x<\tfrac12$, the two coinciding at $\tfrac12$
(\path{fullReliefMinDrift_eq_reflect}, \path{fullReliefExitHorizon_eq_reflect},
\path{entryBound_lt_exitBound_iff_lt_half}). The partial-adjustment law is reflection-symmetric,
so its comparison above is a corollary. The class is larger:
$\F(x,s)=(x+s)/2+(s-x)^3/10$, clipped off the square, is reflection-symmetric,
adjusts strictly and is affine at no rate (\path{cubicReflectF_ne_affineF},
\path{reflectParams_entryBound_lt_exitBound_iff}). Both statements compare two
sufficient durations of temporary shocks, not required durations or equilibrium
behavior.

\paragraph{Why the class under the lifting condition and not corner economies.} \emph{(From the opening of \Cref{sec:four-regimes}.)}

(Several corner economies of
\Cref{sec:existence:nonvacuity} below have the maximal sanction along
equilibrium play. A corner on the equilibrium path does not make the
sanction insensitive to compliance deviations; what the class under the lifting condition adds
is an interior sanction at non-compliance and relief at full compliance,
which is what lets the two forward-looking incentives be separated.) 

\paragraph{The target-only case: the rate, and the sender's state--sanction cross term.} \emph{(From after \Cref{prop:target-only}.)}

On the quadratic--linear family under interiority, which keeps the vertex
$(u_0+rx)/q$ below full sanction so that the myopic sanction rule is the vertex (without
interiority the rule is the vertex clipped at full sanction), the rate in (iii)
is a centered error bound toward the steady state $u_0/(q-r)$, obtained because
the vertex rule is affine about it with slope $r/q$ and the lower
bound on the transition's sanction slope moves the state at least the fraction $d_{\mathrm{lo}}$ of the way; it
is not a pairwise contraction of the closed loop
(\Cref{app:formalization}). The vertex $(u_0+rx)/q$, the rule under
interiority, rises strictly at every positive state with the sender's state--sanction
cross term $r$, and so does the steady state
$u_0/(q-r)$ when the economy with the larger cross term satisfies interiority, which
keeps $r<q$ there (\Cref{app:formalization}): in the
interest-group interpretation of \citet{Kaempfer1988}, a sender whose pro-sanction
constituency responds more to the target's entrenchment imposes more, and
entrenches the target further when the target maintains non-compliance. The comparison is across
economies that share $u_0$ and $q$; the analogue when the sender is also
forward-looking is not a theorem, because the constructed sender value would have to be
pinned down across economies. 

\paragraph{Continuity of the sanction rule.} \emph{(From after \Cref{cor:order-selection-identified}.)}

Continuity of $R_E$ is the mathematical requirement, and it is what
\Cref{prop:order-selection} assumes. One economic condition that supplies it is
uniqueness of the maximizer of $J_S$ at every physical
state and compliance level: a single-valued argmax has a continuous selection, and
the equilibrium already carries continuity of the target policy. We state the
proposition with continuity directly so that the requirement is not confused with
that particular sufficient condition.

\paragraph{Amplification is not claimed.} \emph{(From \Cref{sec:mechanism}.)}

Whether they \emph{amplify} one another is a separate question.
\Cref{thm:four-cells} settles its arithmetic form on the class under the lifting condition: the
common-state cross-difference of the four sanction rules is
strictly positive (\path{four_cell_rule_difference_at_state'},
\Cref{sec:four-regimes}), but that is an identity of the constructed
rules, not strategic complementarity, and no result below signs an
interaction primitive.

\paragraph{Off the corner, in the discussion.} \emph{(From \Cref{sec:discussion:scope}.)}

Off the corner, the Schur condition yields existence and uniqueness within
specified invariant sets of value functions and an atomless transition component yields a purified
equilibrium under stated regularity, while the deterministic joint Bellman
operator need not contract or preserve order; deterministic off-corner
existence is therefore established at example economies rather than
for the full primitive class (\Cref{app:scope-notes}).

\subsection{Scope notes: off-corner existence and the open primitive problem}
\label{app:scope-notes}

The two passages below restrict conclusions that a reader of
\Cref{sec:discussion:scope} might otherwise draw.

The distinction matters most off the corner. The Schur condition yields
existence and uniqueness within specified invariant sets of value functions, and an atomless
transition component yields a purified off-corner equilibrium under the
stated regularity conditions. The deterministic joint Bellman operator,
however, need not be a contraction or preserve order. This is consistent
with the non-existence examples for stationary feedback Stackelberg
equilibrium in finite-state games documented by \citet{Bucarey2023}.
Accordingly, deterministic off-corner existence is established at
example economies rather than for the full primitive class.

A strict ratchet is needed neither for the political state to settle nor for the
limit to be shown to be an extremal steady state; order selection delivers both
without it. What a strict condition buys is a \emph{common} long-run outcome for
a whole interval, and as \Cref{sec:results} records, that is not a separate
requirement but the same one written pointwise: for a continuous, non-decreasing
closed loop mapping an interval into itself, a common limit from every start,
uniqueness of the steady state there, and a strictly signed closed loop are
equivalent. The monotonicity and continuity are not decorative: $G(x)=1-x$ has a
unique fixed point on $[0,1]$ and the sign pattern, and almost every orbit
alternates (\path{nonunique_fixedPoint_breaks_named_limit},
\path{monotone_orbit_tendsto_of_unique}).

Uniqueness inside the claimed interval is necessary. Local isolation of a
candidate limit is insufficient because a second steady state elsewhere retains
its own path. \Cref{prop:derived-interval} obtains uniqueness by ending the
interval at the first steady state. The remaining question concerns an interval
fixed before the equilibrium is solved.

For a pre-specified state interval (the whole high basin, for example),
monotonicity of the sanction rule and the weak level condition do not exclude a
second steady state. A primitive theorem covering such an interval would need to
construct the joint equilibrium, derive a sanction level floor, and establish the
direction of the sanction rule after substituting the target's endogenous
policy. Monotone-argmax
results can supply direction but not level
\citep{MilgromShannon1994,QuahStrulovici2009}; a monotone decision rule also does
not imply that the realized sanction at $x$ is at least $x$, which is the
condition needed for an upward state path. The corner proof cannot simply be
relaxed, because the condition that removes cross-value feedback is also what pins
the sender's response at the ceiling. What the equivalence adds is a target:
any sufficient condition must imply uniqueness of the basin steady state. The
open primitive problem is therefore to derive that uniqueness from payoff and
transition primitives on a pre-specified interval.

\paragraph{Four scope examples.}
\emph{Tied stage maximizers.} With stage payoff $(3/2)x-\tfrac12-s^2(1-s)^2$,
sender value $2y$, the affine transition at rate $\tfrac12$ and $\beta_S=\tfrac12$,
the stage maximizers at state $0$ are $\{0,1\}$, the sender's continuation value at the successor state is
$\Gamma(s)=s$, the stage loss measured from the stage maximizer $0$ is outweighed at $s'=\tfrac12$ by the discounted continuation gain
($\tfrac1{16}<\tfrac14$), and the forward-looking objective
$-s^2(1-s)^2+s/2$ has the unique maximizer $1$: a comparison measured from a non-greatest stage
maximizer orders only that maximizer (\path{tied_stage_maximizers_counterexample}).

\emph{An endpoint floor is not a whole-basin corner.} The rule $R(x)=(x+1)/2$
satisfies $R(x)\ge x$ and $R(1)=1$ with $R(x)<1$ below the top, and is realized
by an equilibrium with $U_S=-(s-(x+1)/2)^2$, $U_T=-a^2$, both values zero and
$\alpha\equiv0$; a floor reaching the top state forces $R(1)=1$ and nothing more
(\path{endpoint_floor_not_whole_basin_corner}: the equilibrium is constructed,
both Bellman equations verified, and the sanction rule shown to be the unique
maximizer of the sender's objective).

\emph{Full compliance in the sender-only case} and \emph{the failed
entrenchment-condition converse}: the two economies in the paragraph after
\Cref{prop:sender-only} (\path{exists_sender_only_full_compliance},
\path{loss_offset_converse_counterexample}; coefficients and values in
\Cref{app:formalization}). A third economy shows that the lifting condition at equality
is also compatible with strictly \emph{partial} compliance: with
$(q,u_0,k,r,w,\ell,c,n,v)=(3,1,2,\tfrac12,1,1,4,0,30)$, the affine transition at
rate $\tfrac12$ and both discounts $\tfrac12$, $A=2=k$ and the zero level is
$a^\ast(x)=\tfrac23+\tfrac x4$ at every state
(\path{exists_equality_placability_partial_compliance}). Equality in
the lifting condition $A\le k$ therefore characterizes nothing about the level.

\subsection{Further exposition of the results}
\label{app:moved-ah}

None of the passages below is a theorem
statement or a hypothesis of a main-text result; each qualifies or
interprets one. Each is headed by what it is and where it stood. The
stage-complementarity route, which is a hypothesis block, moved on the
author's word to \Cref{sec:flip}.

\paragraph{The sender's Lipschitz constant and the transition-adjusted slope condition, in full.} \emph{(From \Cref{sec:assumptions}; the display \eqref{eq:F-adjusted-sender-chord} stays in the main text and is not repeated with its label here.)}

Let $L_{S,x}^{\alpha}$ bound the inherited-state variation of the
sender's \emph{composite} stage payoff along a continuous target policy:
\[
\left|\Us(x_2,\alpha(x_2),s)-\Us(x_1,\alpha(x_1),s)\right|
\le L_{S,x}^{\alpha}|x_2-x_1|.
\]
Under the optional transition regularity \eqref{eq:F-state-lip}, denote
the sender value's envelope Lipschitz bound by $K_S$:
\begin{equation}
K_S=\frac{L_{S,x}^{\alpha}}{1-\beta_S\gamma_F}.
\label{eq:KS-closed}
\end{equation}
This is the standard value-function Lipschitz recursion
\citep{StokeyLucas1989,BertsekasTsitsiklis1996}. The corresponding
target bound uses the inherited-state modulus of the target's
\emph{closed-loop} successor map, which need not equal $\gamma_F$
because the sender response also moves with the state. The target
constant, selector moduli, curvature constants, and the distinction
between primitives and automatically derived intermediates are collected
in \Cref{app:closedforms}. Under the partial-adjustment law,
\eqref{eq:KS-closed} specializes to
$L_{S,x}^{\alpha}/[1-\beta_S(1-\delta)]$.

The modulus $\gamma_F$ does not bound the transition's response to the
sanction. For sanction comparisons the exact general-$\F$ continuation
allowance is instead
\begin{equation}
\beta_S K_S\bigl|\F(x,s_2)-\F(x,s_1)\bigr|.
\label{eq:F-continuation-allowance}
\end{equation}
Accordingly, the cardinal general-$\F$ discharger used below is the
transition-adjusted slope condition
\begin{equation}
\Us(x,a,s_2)-\Us(x,a,s_1)
>
\beta_S K_S\bigl|\F(x,s_2)-\F(x,s_1)\bigr|,
\qquad s_1<s_2,
\end{equation}
on the range used below. It requires no
separate Lipschitz modulus of $\F$ in its sanction argument. Under the
partial-adjustment law, a stage lower slope $m_S$ makes
\eqref{eq:F-adjusted-sender-chord} follow from
$\beta_S K_S\delta<m_S$.

\paragraph{What \textup{(F6-Q)} is and is not.} \emph{(From \Cref{sec:model}, after \eqref{eq:F-state-lip}.)}

Condition \textup{(F6-Q)} is not part of the qualitative baseline
\textup{(F1)--(F5)}. It is stated on the feasible square, and that is the
form the formalization consumes (\Cref{app:formalization}); the global extension of the transition described
above preserves a square-stated modulus on the square but not outside it, so
\textup{(F6-Q)} is not to be read as a global Lipschitz bound. It is the quantitative statement that, holding
today's sanction fixed, inherited political differences are transmitted
at most at rate $\gamma_F$; discounting must dominate that persistence.
Under \eqref{eq:affine}, $\gamma_F=1-\delta$. Thus $\gamma_F$ is the
inherited-state persistence coefficient, not the sanction pass-through
coefficient $\delta$.

\paragraph{Why the derived interval has one steady state, and when it is degenerate.} \emph{(From \Cref{sec:results}, after \Cref{prop:derived-interval}.)}

Uniqueness on $[\bar x,x_s]$ holds by construction because no state below the
least steady state can itself be steady. The interval may nevertheless be
degenerate: if $G_E(\bar x)=\bar x$, then $x_s=\bar x$. Positive length requires
strict motion at the pivot. The example economy in
\Cref{sec:existence:nonvacuity} satisfies this requirement because its sanction
rule equals one against a pivot of $19/20$.

\paragraph{What the three long-run results say about sanctions regimes, in full.} \emph{(From \Cref{sec:results}, the paragraph of that name.)}

\Cref{prop:order-selection} says the inherited political state decides which such
regime the economy reaches, and decides it at date zero: if the equilibrium
sanction at the inherited state already exceeds that state, entrenchment builds
and the economy settles at the \emph{lowest} self-sustaining regime above where it
began; if the sanction falls short, entrenchment erodes to the \emph{highest}
regime below. Where several regimes are reachable, history selects among them,
and it selects within a single equilibrium: no change of expectations and no
coordination on a different continuation is involved. The mechanism is path
dependence, not indeterminacy. Whether several are reachable is a property of the
economy (\Cref{cor:initial-condition-selection} exhibits one where they are),
and where only one is, the proposition says every start reaches it.

\Cref{prop:derived-interval} says more on a narrower range. From any state up to
the first self-sustaining regime above the pivot, every path reaches that regime.
History selects across regimes but not within an interval containing only one.

What none of the three delivers is a rate. They say where a sanctions regime
ends up, not how long it takes to get there, and nothing in the argument bounds
the transition. A regime that is self-sustaining in the limit may be reached only
asymptotically, so the results speak to durability rather than to the period over
which relief becomes implausible.

\paragraph{The derived dividing line: what it determines and the three kinds of threshold.} \emph{(From \Cref{sec:mechanism:bistability}, after \Cref{prop:derived-dividing-line}.)}

The two convergence regions therefore extend strictly on both sides:
from $[0,\bar x_L]$ out to $[0,x_s^-)$, and from $[\bar x_H,1]$ down to
$(x_s^+,1]$. Their endpoints are derived from the equilibrium. The interval
$[x_s^-,x_s^+]$ need not collapse to a point: a continuous map can fix a whole
subinterval, and multiple steady states rule out a common limit across it. Whether
the two bounds coincide in the dual-basin economy remains
undetermined. Continuity is essential for obtaining a middle steady state and is
supplied there by uniqueness of the sender's dynamic argmax.

The notation distinguishes three objects. The stage-payoff tipping states
$\bar x_T$ and $\bar x_S$ are infima of payoff flip sets. The basin boundaries
$\bar x_L$ and $\bar x_H$ are inputs to the conditions for persistence and for convergence to zero.
Only $x_s^-$ and $x_s^+$ are steady states derived from the closed loop.

\paragraph{Strict comparison at a fixed target policy: what it needs and what it does not say.} \emph{(From \Cref{sec:compstat}, after \Cref{prop:strict-dominance}.)}

No concavity, differentiability, uniqueness or positive discount enters, and
$\Gamma$ need not be monotone: one comparison per lower sanction, $\Gamma(d)\le\Gamma(m)$ for
every feasible $d\le m$, suffices
(\path{all_stage_lt_dynamic_of_greatest_stage_improvement_of_wedge_le}), a
non-decreasing $\Gamma$ being one way to supply it; the strict increase of $\Gamma$ that
the first formalization assumed is in any case derived from primitives
(sanctions build entrenchment and the sender values it). The stage loss must be measured from the greatest
stage maximizer: measured from a lower tied one it orders only that one,
and an economy with stage maximizers $\{0,1\}$, the stage loss measured from $0$, and the
maximizer $1$ of $J_S$ shows the all-maximizer conclusion fails
(\Cref{app:formalization}). The proposition orders the two sets of
maximizers and says nothing about interiority: a strict gap is consistent
with both sanctions at corners. It is discharged at the example
economies as the property that no maximizer of $J_S$ is
stage-optimal (\Cref{app:formalization}).

\paragraph{The pathwise ordering: the scope of the corollary.} \emph{(From \Cref{sec:compstat}, after \Cref{cor:pathwise-dominance}.)}

The corollary is stated at the scope of the certificate theorem that proves
it (\Cref{app:formalization}): the statewise inequality comes from the
certificate's myopic ceiling against the sender's floor, not from the
compliance ordering $\alpha\le\alphaM$. It is the corollary's display at the
certificate's own policies; it does not by itself compare
sanctions at two different, endogenously evolving states, and the path
conclusion uses that both paths stay in the basin, the myopic one by
\eqref{eq:myopic-basin-preservation} and the forward-looking one by the
certificate's floor $R_E(x)\ge\max\{x,\bar s\}$ (Step~1 of \Cref{app:proof-pathwise}). The extra hypothesis is substantive because the
corollary compares different state paths. The Eaton--Engers corner economy, where the myopic rule does not preserve a
high basin and the pathwise ordering holds for a simpler reason, is in
\Cref{app:moved-ad}.

\paragraph{Adjustment speed and sender patience in the comparative statics of lock-in.} \emph{(From \Cref{sec:compstat:lockin}.)}

Adjustment speed has a different role. For any equilibrium-derived
sanction rule $\sigma_{\mathrm{real}}$, strict adjustment gives
\[
\F\bigl(x,\sigma_{\mathrm{real}}(x)\bigr)=x
\quad\Longleftrightarrow\quad
\sigma_{\mathrm{real}}(x)=x.
\]
Holding the rule fixed, two positive affine adjustment rates therefore
have the same steady-state set. A higher rate changes the transition
path and convergence time, but it changes a steady state only through
its effect on equilibrium policies or equilibrium selection.

At a fixed continuation value and target policy, greater sender patience
weakly raises the dynamic best response whenever the sender value is
non-decreasing over reachable successor states. The corner economies also
permit a total equilibrium comparison. There, policies remain fixed,
the sender's absorbing-state payoff is positive, so its value is
$U_S(1,0,1)/(1-\beta_S)$. Greater sender patience therefore strictly
raises its lock-in value, while the target value is unchanged because
the sender action remains pinned. \Cref{prop:equilibrium-patience-feedback} gives
the exact direct-plus-continuation decomposition and
\Cref{tab:compstat-windows} reports the parameter intervals over which
the comparison applies.

\paragraph{The constants of the contracting closed loop.} \emph{(From \Cref{sec:compstat:activation-not-hysteresis}, after \Cref{prop:activation-not-hysteresis}.)}

The constants $14/23$ and $35/47$ are evaluated upper bounds,
the values of the proved formula at those members, and not measured
convergence rates or the sharp Lipschitz constants of $G$. The formula rests on the sharp form of the argmax
stability lemma: a $\mu$-strongly concave objective moves its maximizer at
most $L/\mu$ under an $L$-Lipschitz perturbation; the factor-four form of the
first formalization gave $43/46$ at the affine member and $277/235$ at the
non-affine one, where it was inconclusive.

\paragraph{``Finite'' means long enough for the economy, not few.} \emph{(From \Cref{sec:counterfactual}, after \Cref{thm:shock}.)}

Nor is the duration small in any uniform sense: it is
a sufficient bound that depends on the economy through $\eta_{\F}(\bar x)$,
and for any preassigned number of periods there is a partial-adjustment
law slow enough that a temporary shock of that length ends below the
threshold. ``Finite'' means long enough for the economy, not few.

\paragraph{The relief duration: sufficiency, institutional designs, and the two durations.} \emph{(From \Cref{sec:relief-exit}, after \Cref{thm:relief-exit}.)}

The bound $\widehat N^{R}_{\F}(\bar x)$ is a sufficient duration for a
temporary suspension of sanctions. It is not a minimal duration, and
the theorem does not say that the equilibrium would choose such a shock.
Sunset clauses, mandatory review dates, and explicit termination
requirements are distinct institutional designs; the model speaks to them
only insofar as a design commits to a zero-sanction path of the stated
length, which none of them does automatically, and it does not establish
their optimality or their effects on the governments' incentives. The entry
and exit durations are computed from different sections of the same transition
and need not coincide; their comparison, with its specialization to the partial-adjustment law, is in
\Cref{app:moved-ad}.

\paragraph{Two limits of the post-relief statement.} \emph{(From \Cref{sec:relief-exit}, after \Cref{prop:post-relief-persistence}.)}

A second limit is at
the threshold itself: relief
delivers $x\le\bar x$ while the high region of \Cref{thm:shock} is
$[\bar x,1]$, so the two regions share the point $\bar x$, and reaching it
does not by itself select a distinct low outcome; a same-equilibrium
contrast of two such outcomes would need disjoint basins or strict entry into a
low attraction region, which this section does not claim. And the duration
is conservative, being built from a worst-case drift over the whole basin
rather than along the realized path.

\paragraph{What the entry theorem and its corollary do and do not say.} \emph{(From \Cref{sec:relief-exit}, before \Cref{prop:two-sided-one-equilibrium}; it refers to \Cref{thm:shock} and \Cref{cor:shock-transition-strength}.)}

The theorem combines exogenous entry with endogenous persistence. It
does not say that equilibrium sanctions move every initial state into
the high basin. It says that once a temporary sanction shock carries the state across the
threshold, the equilibrium mechanism prevents the high-sanction regime
from unwinding. The corollary compares transition technologies under the
same exogenously set sanction; it is not a comparative static of the equilibrium
sanction itself.

\paragraph{Asymmetric adjustment, inaction bands and path dependence.} \emph{(From \Cref{sec:model}, the primitives.)}

State-dependent and asymmetric adjustment rates remain allowed when
they satisfy \textup{(F1)--(F5)}. Those conditions exclude an inaction
band with $\F(x,s)=x$ when $s\ne x$; under a fixed sanction $s$, the unique
rest point of $\F(\cdot,s)$ is $x=s$. Dual-basin path dependence instead requires the equilibrium sanction
rule to cross the diagonal more than once. The constructions below
exhibit invariant directional basins; they do not establish multiple
crossings as a general fact.
Compliance affects the state only through its effect on the sender's
equilibrium sanction. Asymmetric adjustment connects this mechanism to
political path dependence \citep{Pierson2000} without imposing an
inaction region.

\paragraph{A1: direction, and why direction does not determine level.} \emph{(From \Cref{sec:assumptions}, after \Cref{ass:sender}.)}

A1 is ordinal: after the endogenous target policy is substituted, it
makes the sender's dynamic best response non-decreasing in the state.
Single crossing is a stronger sufficient condition
\citep{MilgromShannon1994}; interval dominance also covers the
non-monotone-continuation case in \Cref{rem:ido}. The premise of
\Cref{thm:dominance} is different: it orders the sender value over
successor states while holding the target policy fixed.

Direction does not determine level. To keep the sender response above
$\max\{x,\bar x\}$, the stage gain from raising sanctions must dominate
the largest continuation loss allowed by the value set. This is exactly
the role of \eqref{eq:F-adjusted-sender-chord}; under the partial-adjustment law it
is implied by $\beta_SK_S\delta<m_S$. The proof of the high-basin theorem keeps the
direction and level arguments separate
(\Cref{app:proofs,app:dischargers}).

\paragraph{A3: what stage concavity does and does not deliver.} \emph{(From \Cref{sec:assumptions}, after \Cref{ass:concavity}.)}

A3 imposes diminishing returns to sanctions, as in standard
public-choice models \citep{Kaempfer1988,Eaton1992}. It makes the myopic
best response unique. It does not ensure uniqueness of the maximizer of
$J_S$ because the continuation term can offset stage curvature.
The smooth existence route therefore imposes curvature on the complete
dynamic objective where needed (\Cref{app:closedforms}). A3 restricts
the change in the stage slope, not its sign.

A3 is not what carries the substitution between compliance and the sanction: decreasing differences of $\Us$
in $(s,a)$ order the greatest stage maximizer with no concavity and no
uniqueness (\Cref{app:moved-ad}).

\paragraph{The increasing-differences condition, and why both forward-looking incentives change behavior together.} \emph{(From \Cref{sec:four-regimes}, after \Cref{prop:placable-family}.)}

The increasing-differences condition says that the stage complementarity between the state and
the sanction covers two forward-looking forces against a monotone
response: saturation of the continuation, and cross-effectiveness of the
transition. Under the partial-adjustment law, and for every additively separable
transition with convex state sections, both vanish, and with a sender payoff
convex in the state the condition reads $0\le r_{\mathrm{lo}}$. The quadratic--linear
family is inhabited at an affine member and at a member whose
transition is not affine at any rate; at both the condition is empty. Since
$\tau_x$, $M_{\mathrm{cx}}$ and $M_{\mathrm{mix}}$ are data of the two-sided
extension of the class (\Cref{app:proofs:placable}, panel~A0), the condition, and
every clause that assumes it, is stated for economies in that extension. The
target's forward-looking incentive changes behavior here \emph{because} the sender is responsive: the
substitution rate covers the sender's whole dynamic marginal return to sanctioning, the sanction
responds to compliance, and both forward-looking incentives change behavior together. The
entrenchment condition does double duty, making the forward-looking target
maintain non-compliance and making a sanction increase induced by the state--sanction complementarity a net gain for the
target's value slope.

\paragraph{The sender-only case: the per-unit relief rate and the forward-looking relief condition.} \emph{(From \Cref{sec:four-regimes}, after \Cref{prop:sender-only}.)}

The forward-looking sender is more demanding than the myopic one (it
still sanctions at $a^{00}(x)$, because of the foresight premium), so the
myopic target pays more for relief; but it still buys it. Below
$a^{01}$, at distance $\Delta=a^{01}-a$, the sanction is at least
$k_{\mathrm{lo}}\Delta/M_{\mathrm{tot}}$, where
$M_{\mathrm{tot}}=\bar Q+\beta_S(\sigma_Cd_{\mathrm{hi}}^2+K_SM_{\mathrm{cs}})$ bounds
the curvature of the sender's dynamic objective in the sanction, the stage
curvature ceiling $\bar Q$ plus the discounted continuation curvature through the
transition's moduli (\Cref{app:proofs:placable}). This bound is a per-unit relief
rate, from the semiconvexity of the sender's objective read as a quadratic
support at its maximizer, and
the forward-looking relief condition (RL), $\ell_{\mathrm{hi}}M_{\mathrm{tot}}<k_{\mathrm{lo}}p_{\mathrm{lo}}$
($\ell M_{\mathrm{tot}}<k(c-n)$ on the quadratic--linear family), says that the
sanction pain avoided per unit of compliance, at the lowest price $p_{\mathrm{lo}}$,
exceeds the unit compliance cost: a strict inequality among the class's
rates and nothing else.

\paragraph{The high-basin theorem: how the proof is organized and what each hypothesis does.} \emph{(From \Cref{sec:results}, after \Cref{thm:persistence}.)}

The proof separates the direction and level of the sender's response.
A1 makes the dynamic best response non-decreasing in the inherited state
after the target policy has been substituted. The transition-adjusted
slope condition in P5 prevents that response from falling below the basin floor.
Together these conditions imply $s_t^E\ge x_t^E$; transition
monotonicity then gives $x_{t+1}^E\ge x_t^E$ and preserves the basin.
Under the partial-adjustment law, a familiar sufficient form of P5 is
$\beta_SK_S\delta<m_S$, where $m_S$ is a lower stage-payoff slope. A local form
with no value sets, and an economy showing that the package selects no limit, are
in \Cref{app:moved-ad}.

The target combined slope condition P4 is logically separate. It yields conclusion (vi),
which verifies that the target's complete objective favors higher
sanctions along the relevant range. Conclusions (i)--(iii) continue to
hold under P1--P3 and P5 alone, and A1 is read only for (iii): (i), (ii)
and (iv) hold under P1, P2 and P5. Conclusions (iv) and (v) record that both
realized actions remain equilibrium-optimal. The proof is in
\Cref{app:proofs}.

The same equilibrium supplies these conclusions for every initial state
in the basin; its construction does not depend on $x_0$. The theorem is
conditional on the stated equilibrium and the package P1--P5.
\Cref{thm:corner-existence} and \Cref{prop:primitive-persistence}
provide unconditional applications under narrower primitive
conditions.

\paragraph{The long-run lock-in level: the argument in brief.} \emph{(From \Cref{sec:results}, after \Cref{prop:longrun}.)}

The state sequence is increasing and bounded. Compactness and continuity
then make every convergent sanction subsequence optimal at the state
limit and consistent with its transition. Strict adjustment implies
that any such sanction limit equals the state limit, so the entire
sanction sequence converges there. Under the partial-adjustment law this identity is
immediate from
\[
s^E_t=\frac{x^E_{t+1}-(1-\delta)x^E_t}{\delta}
\longrightarrow x^\ast.
\]

\paragraph{Entry and exit as opposite-direction crossings.} \emph{(From the opening of \Cref{sec:relief-exit}.)}

The entry result has a counterpart on the exit side. The two are
opposite-direction threshold crossings under temporary shocks; whether one
equilibrium sustains both basins is a further question that neither of
them settles in general, and \Cref{prop:two-sided-one-equilibrium} at the
end of this section is the one economy at which it is settled.
\Cref{thm:shock} spends only the ascent half of \textup{(F5)}, in which sanctions
above the inherited state strictly raise it, and only at the full sanction
below the basin. The exit side spends only the descent half, in which sanctions
below the inherited state strictly lower it, and only at the zero sanction
inside the basin.

\paragraph{Three limits of the two-sided statement.} \emph{(From \Cref{sec:relief-exit}, after \Cref{prop:two-sided-one-equilibrium}.)}

This is the same-equilibrium two-sided statement that
\Cref{sec:relief-exit} declines to make in general: at this economy one
equilibrium sustains both basins, and temporary shocks move between them.
Three limits travel with it. The transition is affine. The two basins do
not exhaust the state space, and nothing is said about the interval between
them. And the statement is an existence certificate: it shows the two level
conditions of the general composition can hold at one equilibrium, and
nothing about other economies.

\paragraph{Predictions inside a high basin.} \emph{(From \Cref{sec:discussion}, predictions and limitations.)}

Inside a high basin satisfying the theorem's conditions, the model predicts non-decreasing
realized sanctions, a weakly positive foresight premium of the sender, relative to a myopic
sender at a common target policy, when the sender value is non-decreasing,
and persistence after a sufficiently long temporary sanction
shock. In the dual-basin construction, initial political conditions
select between convergence to zero under low sanctions and lock-in under high sanctions. Faster
state adjustment shortens the entry duration under the partial-adjustment law and changes the
transition path; holding the sanction rule fixed, it does not
change the steady-state set.

\paragraph{The conclusion's two paragraphs on policy paths and the long run.} \emph{(From \Cref{sec:conclusion}.)}

The mechanism has distinct implications for policy paths. A
forward-looking sender sanctions at least as heavily as a myopic sender
at the same target policy when the political state raises its value. A
temporary sanction shock can move the economy into the high
basin, after which equilibrium behavior keeps the state there. When low
and high basins coexist under one equilibrium, initial political
conditions select between convergence to zero and lock-in.

The long-run results also determine where a monotone equilibrium path
settles. Its initial direction selects the nearest steady state on that
side, and every start between the basin pivot and the first steady state
above it converges to that same state. This common limit is derived from
the equilibrium rule; it need not be interior, and convergence need not
have a uniform rate.
